En la Tab. \ref{tab:data} se muestra el voltaje $V_\text{out}$ nominal y medido, al igual que el
número binario correspondiente. Además, se incluye el error relativo porcentual de
cada medición.

\begin{table}[H]
	\centering
	\caption{Medidas de voltaje}
	\begin{tabular}{cccc}
		\multicolumn{2}{c}{Números} & \multicolumn{2}{c}{Medida} \\
		\hline 
		Binario & Decimal (\unit{\volt}) & $V_\text{out}$ (\unit{\volt}) & $\varepsilon$ (\%) \\
		000 & 0 & $0.0000$ & Na \\
		001 & 1 & $0.996$ & $0.40$ \\
		010 & 2 & $2.001$ & $0.05$ \\
		011 & 3 & $2.992$ & $0.0027$ \\
		100 & 4 & $4.00$ & $0.0$ \\
		101 & 5 & $4.99$ & $0.2$ \\
		110 & 6 & $6.00$ & $0.0$ \\
		111 & 7 & $6.98$ & $0.3$ \\
		\hline 
	\end{tabular}
	\label{tab:data}
\end{table}

Por medio de los casos binarios 001, 010 y 100 se consguió validar los equivalente de thevenin, que correspondian con 1, 2 y 4 voltios respectivamente, el error de estos casos tuvo un error maximo de 0.40\%.

Por medio del resto de casos se valida el modelo de superposición, los cuales tienen un error maximo del 0.3\%.

Los errores se pueden explicar por diferentes factores, como desviaciones de los valores de las resistencias, las cuales tenian una incertidumbre del 5\%, lo que corresponde a una variación maxima de 50 \unit{\ohm}. Otra fuente de error es la estabilidad de la fuente, la cual tenia dificultades para fijar el valor el 16 \unit{\volt} por medio de la perilla además oscilar alrededor del valor requerido.

En terminos de la incertidumbre de las medidas la unica fuente es la propia del instrumento de medición, la cual vario de 0.001 \unit{\volt} para valores menores a 4 \unit{\volt} y 0.01 \unit{\volt} para valores mayores o iguales a 4 \unit{\volt}.
